\section{The Goal Is Committed Upstream, Not Diffed at the Pivot}
\label{sec:mechanism}

We have shown that a decodable, steerable goal encoding is present at the pivot. A natural mechanistic
hypothesis---which we initially pursued---is that the model \emph{computes} this goal at the pivot as a
running \emph{difference} between the task specification and its own progress (a ``\diffmodel'':
$g_k \approx \goal - \progress$). This section shows that account is \textbf{not} what happens. Three
results converge on a different picture: the next goal is \emph{pre-committed upstream} and merely
\emph{read out and routed} by the \goalheads.

\subsection{The goal heads are sufficient but not necessary}
\label{sec:mechanism:ablation}

We test the \goalheads' causal role with two matched ablations on the locked corpus ($n{=}253$; setup in
\Cref{fig:sink-setup}, details in \Cref{app:goalscope}).

\paragraph{Sufficiency (blind-from-span).} We blind \emph{every} upper-layer head from the specification
span \emph{except} the \goalheads. With only the \goalheads able to read what the task asks for, the
model still implements the correct frontier sub-task (near baseline), whereas sparing a random
equal-size, depth-matched set collapses to $\sim$0.09. The \goalheads are \emph{sufficient} to carry the
goal into the implementation.

\paragraph{Necessity (sink-ablation).} We instead silence the \goalheads themselves (restricting them to
attend only to the sink position from the pivot onward). This hurts \emph{no more than} silencing a
random equal-size set: the computation reroutes. The \goalheads are \emph{not necessary}---the goal is
stored redundantly and can be recovered without them.

\begin{figure}[t]
\centering
\begin{subfigure}{0.46\linewidth}
\includegraphics[width=\linewidth]{Figures/ablation_sufficiency_setup.png}
\caption{sink / blind-from-span setup}
\end{subfigure}\hfill
\begin{subfigure}{0.46\linewidth}
\includegraphics[width=\linewidth]{Figures/ablation_sufficiency_necessity.png}
\caption{sufficiency (spare goal heads) vs.\ necessity (silence goal heads)}
\end{subfigure}
\caption{The \goalheads are \emph{sufficient} (spare only them and the model still implements the
sub-task) but \emph{not necessary} (silence them and it reroutes, no worse than random). Placeholders.}
\label{fig:sink-setup}
\end{figure}

Sufficient-but-not-necessary means the \goalheads are best read as a \emph{read-out / routing channel}
for a goal that is represented more diffusely, not as the unique site that \emph{computes} it.

\subsection{The goal is linearly present upstream of the goal-head layers}
\label{sec:mechanism:upstream}

If the \goalheads only route the goal, where is it computed? We probe the residual stream directly with
\textbf{OVReader}: for a residual state $h$ taken after layer $L$ at the pivot, OVReader applies the
\goalheads' output--value (OV) circuit to $h$ and reads the result through \goalscope---asking whether
the upcoming sub-goal is \emph{already present in the residual}, in the exact form a goal head would
write it, from layer $L$ onward. The goal becomes linearly readable well \emph{before} the goal-head
layers concentrate (from $\sim$L40; \Cref{fig:ovreader}). The \goalheads therefore inherit a goal that
upstream computation has already made available, consistent with the necessity result.

\begin{figure}[t]
\centering
\begin{subfigure}{0.46\linewidth}
\includegraphics[width=\linewidth]{Figures/ov_reader_skyline_first_impl.png}
\caption{read-out quality vs.\ probe layer $L$}
\end{subfigure}\hfill
\begin{subfigure}{0.46\linewidth}
\includegraphics[width=\linewidth]{Figures/ov_reader_goalscope_bars.png}
\caption{OVReader vs.\ \goalscope}
\end{subfigure}
\caption{OVReader: the goal is linearly present in the residual from $\sim$L40, before the goal-head
band---i.e.\ computed upstream and read out later (placeholders).}
\label{fig:ovreader}
\end{figure}

\subsection{The pivot query already carries the pre-decided goal, not a difference}
\label{sec:mechanism:precommit}

The sharpest test of the \diffmodel account isolates \emph{what the pivot query carries}. We hold the
model's \emph{progress identical} on both sides and diverge only the \emph{order} of the task list, so
the two hypotheses predict different next tokens:

\begin{center}\small
\begin{tabular}{llll}
\toprule
side & task list (order) & prefilled (progress) & native next \\
\midrule
SRC & $[T_1, T_2, \mathbf{T_3}, T_4]$ & $\{T_1, T_2\}$ & $\mathbf{T_3}$ \\
DST & $[T_1, T_2, \mathbf{T_4}, T_3]$ & $\{T_1, T_2\}$ & $\mathbf{T_4}$ \\
\bottomrule
\end{tabular}
\end{center}

Under \textbf{(H-progress / \diffmodel)}, the query encodes only \emph{how far along} the model is
($\{T_1,T_2\}$, identical on both sides), and the next goal would be computed downstream from the task
keys/values---so transplanting SRC's query into DST should leave DST pursuing $T_4$. Under
\textbf{(H-committed)}, the query already carries the \emph{computed next goal} itself, so the transplant
should flip DST to $T_3$. We observe the latter: transplanting the pivot query flips DST toward $T_3$
(\Cref{fig:precommit}). The next goal is \emph{already decided} by the pivot and carried in the query;
it is not (re)derived at the pivot as a difference against progress.

\begin{figure}[t]
\centering
\includegraphics[width=0.7\linewidth]{Figures/q-transfer-diff-vs-prog.png}
\caption{Query transplant with progress held fixed and task order swapped: DST flips to the source's
\emph{pre-decided} next goal $T_3$, not to the progress-implied $T_4$---evidence against a pivot-time
diff (placeholder).}
\label{fig:precommit}
\end{figure}

\subsection{Which goal, not whether it succeeds}
\label{sec:mechanism:competence}

A direction that predicts \emph{success} could merely track model competence rather than \emph{which}
sub-goal is pursued. Our interventions control for this: the steering and transplant tests
(\Cref{sec:steer,sec:mechanism:precommit}) are \emph{which-goal} contrasts among equally-valid orderings
($T_3$ vs.\ $T_4$, $g_\src$ vs.\ $g_\dst$), where every option is correct, so a shift can only mean the
representation encodes intent, not success. We report the failure modes---cases where the goal signal is
weak or a transplant fails to trigger---in \Cref{app:dissociation}, and discuss them next.
\arnab{fold in the kv-transfer failure taxonomy (naming/datatype; patched read-out reverting to
dst/neither).}

\subsection{Summary}
\label{sec:mechanism:summary}

The next instrumental goal is \emph{committed upstream} (linearly present in the residual by $\sim$L40),
\emph{stored redundantly} (silencing the goal heads reroutes), \emph{carried pre-decided in the pivot
query} (not diffed against progress), and \emph{read out and routed} by a sparse set of \goalheads that
are sufficient but not necessary. This replaces the pivot-time \diffmodel picture with a
commit-then-route account of goal pursuit.
